\section*{Capítulo \ref{ch:g6:decomposers-and-soil} --- Os decompositores e o solo}

\begin{solution}{exo:g6:decomposers-and-soil:1}
Grãos minerais (rocha moída), restos orgânicos em todos os estágios
de \omterm{prop:g6:decomposers-and-soil:decomposition}{decomposição}, água e ar --- além da população viva dele. A camada
escura é o \omterm{def:g6:decomposers-and-soil:soil}{húmus}.
\end{solution}

\begin{solution}{exo:g6:decomposers-and-soil:2}
De \omterm{def:g6:plants-are-producers:matter}{matéria orgânica} morta --- \omterm{def:g1:plants-around-us:parts}{folhas} caídas, madeira morta, fezes,
cadáveres. Visíveis: minhocas, tatuzinhos, piolhos-de-cobra,
colêmbolos (ou besouros do esterco, larvas de mosca). O que se
alimenta com um feltro de fios: um \omterm{def:g6:decomposers-and-soil:decomposers}{fungo}.
\end{solution}

\begin{solution}{exo:g6:decomposers-and-soil:3}
De volta em \omterm{def:g6:plants-are-producers:matter}{matéria mineral} --- \omterm{def:g6:decomposers-and-soil:soil}{húmus}, depois \omterm{prop:g4:what-plants-need:needs}{sais minerais}
dissolvidos e gás para o ar. Ela inverte o ofício das produtoras, que
construíram o orgânico a partir do mineral.
\end{solution}

\begin{solution}{exo:g6:decomposers-and-soil:4}
Novembro: caída inteira. Inverno: amolecida e minada pelos \omterm{def:g6:decomposers-and-soil:decomposers}{fungos},
rendada até virar \omterm{def:g3:bones-and-muscles:skeleton}{esqueleto} por colêmbolos e tatuzinhos. Primavera:
arrastada para baixo e comida por uma minhoca, expelida como
migalhas. Verão: \omterm{def:g6:decomposers-and-soil:soil}{húmus} preto; os minerais dela na água do \omterm{def:g6:decomposers-and-soil:soil}{solo},
bebidos pelas \omterm{def:g1:plants-around-us:parts}{raízes} --- talvez as do próprio carvalho.
\end{solution}

\begin{solution}{exo:g6:decomposers-and-soil:5}
Porque os \omterm{def:g6:decomposers-and-soil:decomposers}{decompositores} o comem tão rápido quanto os anos o soltam:
cada tapete é desmontado em \omterm{def:g6:decomposers-and-soil:soil}{húmus} e minerais antes de o seguinte cair
por inteiro.
\end{solution}

\begin{solution}{exo:g6:decomposers-and-soil:6}
Calor, \omterm{def:g6:exploring-our-environment:conditions}{umidade} e ar. O trabalho quase para no inverno --- pelo frio
--- e é por isso que o tapete fica grosso até a primavera quente e
úmida limpar tudo.
\end{solution}

\begin{solution}{exo:g6:decomposers-and-soil:7}
Produtoras $\rightarrow$ \omterm{def:g5:food-webs:roles}{consumidores} (as \omterm{def:g1:plants-around-us:plant}{plantas} comidas);
produtoras e \omterm{def:g5:food-webs:roles}{consumidores} $\rightarrow$ restos mortos, fezes, \omterm{def:g1:plants-around-us:parts}{folhas}
caídas $\rightarrow$ \omterm{def:g6:decomposers-and-soil:decomposers}{decompositores} (os mortos comidos);
\omterm{def:g6:decomposers-and-soil:decomposers}{decompositores} $\rightarrow$ \omterm{prop:g4:what-plants-need:needs}{sais minerais} na água do \omterm{def:g6:decomposers-and-soil:soil}{solo}
$\rightarrow$ bebidos pelas produtoras: uma roda fechada.
\end{solution}

\begin{solution}{exo:g6:decomposers-and-soil:8}
Trabalhadores: minhocas, tatuzinhos, colêmbolos, \omterm{def:g6:decomposers-and-soil:decomposers}{fungos} e as equipes
invisíveis. Matéria-prima: o lixo orgânico da cozinha. Produto: \omterm{def:g6:decomposers-and-soil:soil}{húmus}
escuro, cujos minerais alimentam os canteiros. Ela esquenta porque os
trabalhadores, como todos os \omterm{def:g5:food-webs:roles}{consumidores}, queimam parte do que
comem --- a pilha é aquecida pelo próprio viver da equipe dela.
\end{solution}

\begin{solution}{exo:g6:decomposers-and-soil:9}
Decompositora --- arrastando \omterm{def:g1:plants-around-us:parts}{folhas} para o subsolo e comendo matéria
decomposta --- e aradora: as galerias dela arejam o \omterm{def:g6:decomposers-and-soil:soil}{solo} e os dejetos
dela amassam o mineral e o orgânico juntos. O \omterm{def:g6:decomposers-and-soil:soil}{solo} de campos inteiros
passa pelo intestino das minhocas, e é por isso que os naturalistas
antigos coroaram a minhoca como fabricante da terra de cultivo.
\end{solution}

\begin{solution}{exo:g6:decomposers-and-soil:10}
Canteiro com cobertura morta: \omterm{def:g6:decomposers-and-soil:soil}{húmus} escuro e esfarelado; um censo
rico --- minhocas, tatuzinhos, colêmbolos; a folhagem da superfície
sumindo visivelmente estação após estação; veredito: um ciclo girando,
que vale a \omterm{def:g1:animals-around-us:covering}{pena} alimentar. Trilha pisoteada: poeira pálida e
compactada; uma bandeja quase vazia; o que cai fica sem apodrecer;
veredito: um \omterm{def:g6:decomposers-and-soil:soil}{solo} faminto e sem ar, cuja equipe foi embora.
\end{solution}

\begin{solution}{exo:g6:decomposers-and-soil:11}
O ar. Lacrado e encharcado, o saco sem ar barra os trabalhadores que
precisam de ar e deixa correr só um apodrecimento azedo e sem ar. O
conserto: tirar do saco e pôr numa pilha solta, úmida e revolvida ---
com a \omterm{def:g6:exploring-our-environment:conditions}{umidade} mantida e o ar admitido.
\end{solution}

\begin{solution}{exo:g6:decomposers-and-soil:12}
O que se acumula: os mortos --- os tapetes de todo outono, os restos
de todas as gerações, montanhas de \omterm{def:g6:plants-are-producers:matter}{matéria orgânica} não desmontada. O
que acaba: os \omterm{prop:g4:what-plants-need:needs}{sais minerais} do \omterm{def:g6:decomposers-and-soil:soil}{solo}, nunca reabastecidos depois que a
seta de retorno é cortada; as produtoras bebem os últimos minerais
dissolvidos e depois não conseguem construir mais nada, por mais \omterm{def:g6:exploring-our-environment:conditions}{luz}
que haja. As duas metades são a mesma seta cortada --- a \omterm{prop:g6:decomposers-and-soil:decomposition}{decomposição}
--- vista das duas pontas.
\end{solution}

\begin{solution}{pb:g6:decomposers-and-soil:1}
\textbf{1.} Para que os sacos difiram em uma coisa só --- a malha --- e
qualquer diferença de peso possa ser atribuída a ela: o teste justo,
enterrado.

\textbf{2.} Se a vida invisível do \omterm{def:g6:decomposers-and-soil:soil}{solo} consegue sozinha desmontar
\omterm{def:g1:plants-around-us:parts}{folhas} --- o saco de malha grossa mistura trabalhadores visíveis e
invisíveis; o de malha fina isola os invisíveis.

\textbf{3.} Se as \omterm{def:g6:exploring-our-environment:conditions}{condições} comandam o trabalho: as mesmas \omterm{def:g1:plants-around-us:parts}{folhas}, a
mesma malha, mas cascalho seco contra barranco úmido --- ele testa em
campo a condição de \omterm{def:g6:exploring-our-environment:conditions}{umidade} da
\cref{prop:g6:decomposers-and-soil:speed}.

\textbf{4.} A malha grossa perde mais, a fina menos, o lacrado quase
nada --- previsto direto a partir de quem e do que cada saco admite.

\textbf{5.} Malha grossa: a equipe completa trabalhou --- desmontagem
mais rápida. Malha fina: só os invisíveis trabalharam --- perda real,
mas mais lenta. Lacrado: nenhum trabalhador e nenhum ar --- quase
nenhuma \omterm{prop:g6:decomposers-and-soil:decomposition}{decomposição} de verdade, só um estrago azedo e sem ar.

\textbf{6.} Os fios dos \omterm{def:g6:decomposers-and-soil:decomposers}{fungos} minam e amolecem a \omterm{def:g1:plants-around-us:parts}{folha}; os colêmbolos
e os tatuzinhos a roem até virar renda; a minhoca arrasta os
fragmentos para baixo, come esses fragmentos e expele migalhas
trabalhadas; as equipes invisíveis terminam cada fragmento em \omterm{def:g6:decomposers-and-soil:soil}{húmus} e
minerais.

\textbf{7.} Que ela é uma verdadeira força de trabalho de \omterm{prop:g6:decomposers-and-soil:decomposition}{decomposição}
por si só: o peso saiu do saco com todos os \omterm{def:g1:animals-around-us:animal}{animais} visíveis barrados,
então a desmontagem continuou na escala abaixo da vista.

\textbf{8.} A \omterm{def:g6:exploring-our-environment:conditions}{umidade} acima de tudo --- o cascalho drena e seca --- e,
com ela, a população trabalhadora que um barranco úmido abriga; as
\omterm{def:g6:exploring-our-environment:conditions}{condições} secas deixam ociosa qualquer equipe que chegue.

\textbf{9.} Para o \omterm{def:g6:decomposers-and-soil:soil}{solo} e para o ar: uma parte virou \omterm{def:g6:decomposers-and-soil:soil}{húmus} e \omterm{prop:g4:what-plants-need:needs}{sais
minerais} dissolvidos na água do \omterm{def:g6:decomposers-and-soil:soil}{solo} do barranco, e outra parte foi
queimada pelos \omterm{def:g6:decomposers-and-soil:decomposers}{decompositores} em atividade e saiu como gás --- os
gramas perdidos estão todos contabilizados, nenhum sumiu.

\textbf{10.} Não é coincidência: os minerais das \omterm{def:g1:plants-around-us:parts}{folhas} enterradas
foram liberados exatamente ali, e as urtigas são famosas comilonas de
terreno rico. A faixa verde é a última seta do ciclo tornada visível
--- a \omterm{prop:g6:decomposers-and-soil:decomposition}{decomposição} alimentando as produtoras.

\textbf{11.} O embrulho do aterro é o saco lacrado em escala de
cidade: sem ar e sem equipe, ele nega à \omterm{prop:g6:decomposers-and-soil:decomposition}{decomposição} os trabalhadores
e as \omterm{def:g6:exploring-our-environment:conditions}{condições} dela, então o lixo dura como durou a lama do saco ---
enquanto as mesmas cascas numa pilha arejada, úmida e cheia de
trabalhadores voltam a ser \omterm{def:g6:decomposers-and-soil:soil}{solo} em meses.

\textbf{12.} O inverno provou que o tapete é limpo pelos
\omterm{def:g6:decomposers-and-soil:decomposers}{decompositores} do \omterm{def:g6:decomposers-and-soil:soil}{solo} --- a equipe visível e as equipes invisíveis
juntas, mais rápido onde há calor, \omterm{def:g6:exploring-our-environment:conditions}{umidade} e ar. Essa limpeza
alimenta as reservas minerais do \omterm{def:g6:decomposers-and-soil:soil}{solo} e, por meio delas, as
produtoras: a roda da matéria, girando num barranco de cerca viva.
\end{solution}
